\documentclass[10pt]{extarticle}
\usepackage[utf8]{inputenc}
\usepackage{graphicx}
\usepackage{caption}
\usepackage{amsmath}
\usepackage{amssymb}
\usepackage{amsfonts}
\usepackage[a4paper, margin=1cm]{geometry}
\title{Teoria di \textit{Algebra Lineare}}
\author{Nicolò Giuliani}
\date{28/05/2025}

\begin{document}
\maketitle
\section{Matrici e sistemi lineari}
\subsection{Proprieta' delle matrici}
\begin{itemize}
    \item proprieta' associativa : $(AB)C=A(BC)$
    \item proprieta' distributiva: $A(B+C)=AB+AC$
    \item proprieta' trasposta: $(AB)^T=A^TB^T$
\end{itemize}
Le \textbf{operazioni elementari} su una matrice sono:
\begin{itemize}
    \item $R_i \cdot \alpha, \quad \alpha \in \mathbb{R}\ \{0\}$
    \item $R_i - R_j$
    \item $R_i \overset{\text{Gauss}}{\iff} R_j$
\end{itemize}
\subsection{Risolvere un sistema lineare}
Dato $A\vec{x}=\vec{b}$, con $x=(x_1,\dots,x_n)$:
\begin{itemize}
    \item $\exists! \text{ 1 soluzione } \iff rk(A \mid b)=rk(A)=n$
    \item $\exists \, \infty \text{ soluzioni dipendenti da $n-rk(A)$ parametri (o incognite)} \iff rk(A \mid b) = rk (A) < n$
\end{itemize}

\section{Spazi vettoriali}
E' un insieme ($V$) su un campo $\mathbb{R}$ che gode di queste proprieta': 
\begin{align*}
    \text{somma}: \, V \times V = V \\
    \text{prodotto}: \, \mathbb{R} \times V = V
\end{align*} 
In tal modo porta le seguenti proprieta':
\begin{itemize}
    \item $\vec{u}+\vec{v}=\vec{v}+\vec{u} \quad (\vec{v},\vec{u} \in V)$
    \item $(\vec{v}+\vec{w})+\vec{u}=(\vec{v}+\vec{u})+\vec{w}$
    \item $\exists! \,\vec{0} \iff \vec{0} + \vec{v} = \vec{v}$
    \item $\exists! \, (-\vec{v}) \to \vec{v}+(-\vec{v})=\vec{0}$
    \item $1 \cdot \vec{v}=\vec{v}$
    \item $(\lambda \gamma)\vec{u}=\lambda(\gamma \vec{u}), \quad \lambda,\gamma \in \mathbb{R}$
    \item $\lambda(\vec{v}+\vec{u})=\lambda\vec{v}+\lambda\vec{u}, \quad \lambda \in \mathbb{R}$
    \item $\vec{v}(\lambda+\gamma)=\lambda\vec{v}+\gamma\vec{v}, \quad \lambda,\gamma \in \mathbb{R}$
\end{itemize}
\subsection{Sottospazi vettoriali}
Si dice sottospazio vettoriale $W \subset V$ se:
\begin{align*}
    \vec{v},\vec{u} \in W \to \vec{v}+\vec{u}\in W \\
    \vec{v}\in W \to \lambda\vec{v}\in W, \, \forall \lambda \in \mathbb{R} \\
    \vec{0_v} \in W
\end{align*}

\section{Combinazioni lineari}
Per definizione di spazio vettoriale sappiano che se un vettore appartiene ad uno spazio vettoriale allora anche tutti i suoi multipli lo faranno. \\
Infatti si definisce \textbf{combinazione lineare} il vettore $\vec{v}=\lambda_1 \vec{v_1}+\dots+\lambda_n \vec{v_n} \, \in V$. \\
Quindi se $\vec{v} \in W$ vuol dire che il sottospazio $W$ e' generato dalle combinazioni lineare di $\vec{v}_i$. Ovvero $W=span\{\vec{v_1},\dots,\vec{v_n}\}$, quest'ultimi si dice che \textit{generano} $W$. \\

\noindent\vspace{5px} 
\textbf{Vettore linermente indipendente}: dei vettori $\vec{v_1},\dots,\vec{v_n}$ si dicono linearmente indipendenti solo se 
\begin{align*}
    \alpha\vec{v_1}+\beta\vec{v_2}+\dots+\gamma\vec{v_n}=\vec{0_v} \quad \to \alpha=\beta=\gamma=0
\end{align*}
Dimostrazione della contronominale: 
\begin{align*}
    \text{Ipotizzo: } \vec{v_2}=2\vec{v_1} \neq \vec{0_v}\\
    \Rightarrow \alpha\vec{v_1}+2\beta\vec{v_1}=\vec{0_v} \\
    \Rightarrow \vec{v_1}(\alpha+2\beta)=\vec{0_v} \iff \alpha=-2\beta\\
   \text{Ipotizzo } \beta = 1 \Rightarrow \alpha = -2 \\
    \Rightarrow (-2) \vec{v}_1 + 1 \cdot (2 \vec{v}_1) = \vec{0}
\end{align*}
Quindi esiste una soluzione non banale all’equazione, per cui \(\vec{v}_1, \vec{v}_2\) non sono linearmente indipendenti. 

\section{Basi e dimensioni}
Con base intendiamo un insieme minimo di vettori che generano un insieme $V$.
\begin{align*}
    \beta_{V}=\{\vec{v_1},\dots,\vec{v_n}\} \quad \neq \{\vec{0_v}\}
\end{align*}
Dimostrazione: \\
\vspace{0.5em} Dividiamo la dimostrazione in due parti:
\begin{itemize}
    \item Caso banale: \(v_1,\dots,v_n \quad \neq \vec{0_v}\) e' gia' linearmente indipendente. \\
    Quindi possiamo trovare una base: $\beta=\{\lambda_1\vec{v_1},\dots,\lambda_n\vec{v_n} \mid \lambda_1,\dots,\lambda_n \in \mathbb{R}\}$
    \item Caso induttivo: \(\vec{v_1},\dots,\vec{v_n} \quad \neq \vec{0_v}\) sono linearmente dipendenti; ipotizzo che $\vec{v_n}=\lambda \vec{v}_{n-1}, \, \exists \lambda \in \mathbb{R}$. \\
    Allora $span\{\vec{v_1},\dots,\vec{v}_{n-1}\}=span\{\vec{v_1},\dots,\vec{v}_{n}\}$. Quindi $\exists \beta=\{\lambda_1\vec{v_1},\dots,\lambda_{n-1}\vec{v}_{n-1} \mid \lambda_1,\dots,\lambda_{n-1} \in \mathbb{R}\}$.
\end{itemize}
\textbf{Teorema del completamento}: Per  generare uno spazio $\mathbb{R}^n$ servira' una base composta da precisamente $n$ vettori. Quindi se prendo i vettori linermente indipendenti $\vec{v_1},\dots,\vec{v_k}\in\mathbb{R}^n \, (k<n)$, e una base che GENERA $\mathbb{R}^n$ ovvero: $\beta_n=\{\vec{w_1},\dots,\vec{w_n}\}$; allora $\exists \, \beta_{n}'=\{\vec{v_1},\dots,\vec{v_n},\vec{w_1},\dots,\vec{w}_{n-k}\}$ che genera $\mathbb{R}^n$.\\
\vspace{1em}\\ Quindi per definizione sappiamo che tutte le basi che generano uno spazio vettoriale avranno la stessa dimensione. Sappiamo anche che la dimensione di un sottoinsieme sara sempre $\leq$ del suo spazio vettoriale. \\
\vspace{1em}\\\textbf{Teorema dell'unicità della rappresentazione di un vettore rispetto a una base}: Data la base ordinata generatrice $\beta_{\mathbb{R}^n}=\{\vec{v_1},\dots,\vec{v_n}\}$, allora 
\begin{align*}
    \exists! \text{ la tupla degli scalari } (\alpha_1,\dots,\alpha_n) \iff \vec{v}=\alpha \vec{v_1}+\dots +\alpha_n \vec{v_n} \quad \in \mathbb{R}^n
\end{align*}
Dimostrazione su $\mathbb{R}^2$:
\begin{align*}
    \beta_{\mathbb{R}^2}=\{(1,0),(0,1)\} \\
    \vec{v}=\alpha(1,0)+\gamma(0,1) \quad \in \mathbb{R}^2\\
    \text{Ipotizzo: $\alpha=2,\gamma=3$}\\
    \Rightarrow \vec{v}=(2,3) \\
    \exists! \text{ 2-upla ordinata: } (2,3)
\end{align*}
Precisazione: la $n-upla$ si chiama \textit{cordinate} di $\vec{v}\in \mathbb{R}^n$ sulla base $\beta_{\mathbb{R}^n}$; ovvero $[\vec{v}]_{\beta_V}=(\alpha_1,\dots,\alpha_n)$
\vspace{2em}\subsection{Approfondimento delle matrici}
Una matrice $A \in M_{m,n}(\mathbb{R})=\begin{pmatrix}
    [\vec{v_1}] & \dots & [\vec{v_n}] \\
    \vdots & & \vdots
\end{pmatrix}$ genera un sottospazio di $\mathbb{R}^m$; la sua matrice derivata dall'algoritmo di Gauss genera sempre lo stesso sottospazio. Piu' precisamente i vettori NON NULLI a scala sono linearmente indipendenti, quindi saranno sempre questi in entrambe le due matrici ad generare il sottospazio di $\mathbb{R}^m$. 

\section{Applicazioni Lineari}
\begin{align*}
    f \, : \mathbb{R}^n \to \mathbb{R}^m
\end{align*}
Una applicazioni lineare e' una funzione a $n$ variabili che gode delle seguenti proprieta':
\begin{itemize}
    \item $F(\vec{v}+\vec{u})=F(\vec{v})+F(\vec{u}), \quad \forall \vec{v},\vec{u} \in \mathbb{R}^n$ 
    \item $F(\lambda\vec{v})=\lambda F(\vec{v}), \quad \forall \lambda \in \mathbb{R}, \, \forall \vec{v} \in \mathbb{R}^n $
\end{itemize}
\textbf{Teorema di esistenza e unicità dell'applicazione lineare su una base}: \\
Dati due spazi vettoriale $\mathbb{R}^n$ e $\mathbb{R}^m$ e scegliamo $n$-vettori, anche uguali, di $\mathbb{R}^m$ allora
\begin{align*}
    \exists! \, F(\vec{v}_i)=\vec{w_i} \quad \vec{v_i}\in \mathbb{R}^n, \vec{w_i}\in\mathbb{R}^m \quad \mid \forall i=1,\dots,i=n
\end{align*}
Dimostrazione:
\begin{align*}
    \text{Definiamo: } F(\vec{v})=\lambda_1\vec{w_1}+\dots+\lambda_n\vec{w_n}\\
   \text{Ipotizziamo: } \vec{v}=\lambda_1\vec{v_1}+\dots+\lambda_n\vec{v_n}, \, \vec{u}=\alpha_1\vec{v_1}+\dots+ \alpha_n\vec{v_n} \,(\in \mathbb{R}^n)\\
   \Rightarrow_1 \text{Dimostriamo se e' chiuso rispetto la somma}: \\
   F(\vec{v}+\vec{u})=F((\lambda_1\vec{v_1},\dots,\lambda_n\vec{v_n})+(\alpha_1\vec{v_1},\dots,\alpha_n\vec{v_n})) \\
   =F(\vec{v_1}(\lambda_1+\alpha_1),\dots,\vec{v_n}(\lambda_n+\alpha_n))\\
   =(\lambda_1+\alpha_1)F(\vec{v_1})+\dots+(\lambda_n+\alpha_n)F(\vec{v_n}) \\
   =(\lambda_1+\alpha_1)\vec{w_1}+\dots+(\lambda_n+\alpha_n)\vec{w_n} \quad (\vec{w_1},\dots,\vec{w_n} \in \mathbb{R}^m)\\ 
   \iff \lambda_1\vec{w_1}+\alpha_1\vec{w_1}+\dots+\lambda_n\vec{w_n}+\alpha_n\vec{w_n} \\
   = F(\vec{v}) + F(\vec{u}) \\
    \Rightarrow_2 \text{Dimostriamo se e' chiuso rispetto al prodotto}: \\
    F(\gamma\vec{v})=F(\gamma\lambda_1\vec{v_1},\dots,\gamma\lambda_n \vec{v}_n)\\
    =\gamma\lambda_1F(\vec{v_1})+\dots+\gamma\lambda_nF(\vec{v_n})\\
    = \gamma\lambda_1\vec{w_1}+\dots+\gamma\lambda_n\vec{w_n} \\
    = \gamma (\lambda_1\vec{w_1}+\dots+\lambda_n\vec{w_n}) \\
    = \gamma F(\vec{v})
\end{align*}
\textbf{Teorema di rappresentazione delle applicazioni lineari}: possiamo rappresentare una applicazione linare con \\
1) Come combinazione lineare delle incognite
\begin{align*}
    F(x_1,\dots,x_n)=\begin{cases}
        \alpha_{11} x_{11}+\dots+\alpha_{1n}x_{1n}\\
        \vdots \\
        \alpha_{m1} x_{m1}+\dots+\alpha_{mn}
    \end{cases}
\end{align*}
2) Rispetto ai vettori della base canonica
\begin{align*}
    F(e_1)=\dots \\
    \vdots \\
    F(e_n)=\dots
\end{align*}
3) Come prodotto tra matrice e vettore
\begin{align*}
    F(\vec{x})=A\vec{x} \\
    \begin{pmatrix}
        a_{11} & \dots & a_{1n} \\
        \vdots \\
        a_{m1} & \dots & a_{mn}
    \end{pmatrix}\begin{pmatrix}
        x_1 \\ \vdots \\ x_n
    \end{pmatrix}
\end{align*}
\vspace{1em}
\textbf{Applicazioni lineari composte}: dato $F : \mathbb{R}^n \to \mathbb{R}^m, \, G : \mathbb{R}^z \to \mathbb{R}^n$ allora $\exists! \, F \circ G :  \mathbb{R}^z \to \mathbb{R}^m$. La sua matrice sara' invece $[M]_{F \circ G}=[M]_F \times [M]_G$.
\vspace{2em}
\subsection{Kernel} 
E' lo spazio generato dai vettori che hanno come immagine il vettore nullo.
\begin{align*}
    Kernel=span\{\vec{v} \in \mathbb{R}^n \mid F(\vec{v})={\vec{0_v}}\}
\end{align*}
Dimostrazione:
\begin{align*}
    \text{Dato: } Kernel=span\{\vec{v}+\vec{u}\}\neq\{\vec{0_v}\} \\
    \Rightarrow_1 \text{Dimostriamo se esiste una appl. lin. chiusa rispetto la somma: } \\
    F(\vec{v}+\vec{u})=F(\vec{v})+F(\vec{u})=\{0_v\}\\
    \vec{0_v}+\vec{0_v}=\vec{0_v} \\
    \Rightarrow_2 \text{Dimostriamo se esiste una appl. lin. chiusa rispetto il prodotto: } \\
    F(\lambda\vec{v})=\vec{0_v} \\
    \lambda F(\vec{v})=\vec{0_v} \\
    \lambda \vec{0_v}=\vec{0_v}
\end{align*}
\subsection{Immagine}
E' un sottoinsieme del codominio tale che
\begin{align*}
    Im(F)=span\{\vec{w}\in\mathbb{R}^m \mid \exists \, F(\vec{v})=\vec{w}\} \subseteq \mathbb{R}^m
\end{align*}
Dimostrazione:
\begin{align*}
    \text{Dato: } Im(F)=span\{\vec{c_1}+\vec{c_2}\}\neq\{\vec{0_v}\} \\
    \Rightarrow_1 \text{Dimostriamo se esiste una appl. lin. chiusa rispetto la somma: } \\
    \vec{c_1}+\vec{c_2}=F(\vec{v}_1)+F(\vec{v_2}) \\
    F(\vec{v_1}+\vec{v_2}) \\
    \Rightarrow_2 \text{Dimostriamo se esiste una appl. lin. chiusa rispetto il prodotto: } \\
    \lambda\vec{c_1}=\lambda F(\vec{v_1}) \\
    F(\lambda\vec{v_1})
\end{align*}
\textbf{Proposizione}: L'immagine puo essere generata dall'output della applicazione lineare che ha i vettori in input in diverse basi che generano sempre come spazio vettoriale il dominio.
\begin{align*}
    span\{\beta_{\mathbb{R}^n}\}=Im(F)
\end{align*}
Dimostrazione:
\begin{itemize}
    \item Sappiamo che il dominio puo' essere generato da qualsiasi base di $n$-vettori
    \item Ad ogni base ordinata (che genera il dominio) avra una sua matrice composta dai vettori ordinati
    \item Sappiamo anche che ogni matrice genera un sottospazio del codominio.
    \item Quindi dato che i vettori della base sono per definizione linearmente indipendenti, allora tutti i vettori della applicazione lineare che appartengono al dominio apparteranno a tutte le basi di quella dimensione.
    \item Quindi tutti i vettori generano lo stesso sottospazio nel codominio, quindi l'immagine sara' sempre uguale.
\end{itemize}
\begin{align*}
    \text{Dato: } span\{\beta_{\mathbb{R}^n}\}=\mathbb{R}^n=span\{\beta_{\mathbb{R}^n}'\} \\
    \text{Sappiamo: } \exists \, \vec{v_1},\dots,\vec{v_n}\in\mathbb{R}^n \\
    \text{Sappiamo: } span\{[M]_{\beta_{\mathbb{R}^n}}\}=span\{[M]_{\beta_{\mathbb{R}^n}'}\}\subseteq \mathbb{R}^m \\
    \to \vec{v_1},\dots,\vec{v_n}\in \mathbb{R}^n=span\{\beta_{\mathbb{R}^n}\}\\
    span\{F(\vec{v_i})\}=Im(F) \iff \forall\vec{v_i} \in span\{\beta_{\mathbb{R}^n}^i\} 
\end{align*}
\vspace{1em}\\
\subsection{Caratteristiche delle applicazioni lineari}
\begin{itemize}
    \item Iniettiva: $\forall \vec{v},\vec{u}, \, F(\vec{v})=F(\vec{u}) \iff \vec{v}=\vec{u}$ con $F : \mathbb{R}^n \to \mathbb{R}^m$
    \begin{align*}
        Kernel(F)=\{\vec{0_v}\}
    \end{align*}
Dimostrazione:\\
\begin{align*}
    \text{assumo F iniettiva} \\
    \Rightarrow \quad \forall \, \vec{v}\neq\vec{u} \in \mathbb{R}^n \to F(\vec{v})\neq F(\vec{u}) \\
    [A](\vec{v})\neq[A](\vec{u}) \\
    \text{Sappiamo per definizione: } Kernel= \quad [A](\vec{x})=\{\vec{0_v}\} \\
    \exists! \, \vec{x} \iff A\vec{x}=\{\vec{0_v}\} \\
    A^{-1} \cdot A\vec{x}=\{\vec{0_v}\}\cdot A^{-1} \\
    \to I\vec{x}=\{\vec{0_v}\} \\
    \to \vec{x}=\{\vec{0_v}\} \\
\end{align*}
    \item Suriettiva: $\forall \vec{w}, \, \vec{w}=F(v) \to \exists \vec{v}$ con $F : \mathbb{R}^n \to \mathbb{R}^n$
    \begin{align*}
        Im(F)=\mathbb{R}^m
    \end{align*}
    Dimostrazione:
    \begin{align*}
        \text{assumo F suriettiva} \\
        \forall \vec{w}, \, \vec{w}=F(\vec{v})\\
        \exists \vec{v_i} \to F(v_i)\in\mathbb{R}^m \\
        \mathbb{R}^m=F(v_1)+\dots+F(v_n)\\
    \end{align*}
    \item Invertibile: se biettiva e $\exists F^{-1}(\vec{w})=\vec{v} \iff F(\vec{v})=\vec{w}; \quad F : V \to W$ \\
    \begin{align*}
        \text{Si dicono isomorfi due spazi vettoriali se:} \\
        dim(V)=dim(W) \\
        \to V \equiv W
    \end{align*}
\end{itemize}
Prima abbiamo definito che 
\begin{align*}
    \exists F^{-1}(\vec{w})=\vec{v} \to F(\vec{v})=\vec{w}
\end{align*}
questa si chiama \textit{contro-immagine}. \\
\vspace{1em} \\
Possiamo definire che la contro-immagine esiste se esiste almeno un vettore $\vec{w} \in Im(F)$. \\
\vspace{15em}\\
\subsection{Teorema della dimensione}
\[
\begin{array}{c}
    \dim(\mathbb{R}^n) = \dim(\ker F) + \dim(\operatorname{Im} (F)) \\
    \rotatebox{90}{$\iff$} \\
    \dim(\ker (F)) = \dim(\mathbb{R}^n) - \dim(\operatorname{Im} (F))
\end{array}
\]
Dimostrazione:
\begin{align*}
    n=dim(Ker) \, + \, dim(Im(F)) \\
    \text{Sappiamo: } Ker=span\{\beta_{Ker}\} \\
    \text{Definiamo: } \beta_{Ker}=\{\vec{v_1},\dots,\vec{v_k}\} \quad (k \leq n) \\
    \text{Cerchiamo altri $n-k$ vettori in $\mathbb{R}^n$}: \vec{u_1},\dots,\vec{u}_{n-k} \\
    \text{Sappiamo da definizione di Kernel: } dim(Im(F(\vec{v_i})))=0 \\
    \text{Completiamo la base: } \beta=\{\vec{v_1},\dots,\vec{v_k},\vec{u_1},\dots,\vec{u}_{n-k}\} \\
    \text{Allora: } dim(\beta)=n \\
    \Rightarrow dim(\beta)=dim(Ker) \, + \, dim(Im(F)) \\
    \to dim(\beta)=k-dim(Im(F)) \\ 
    \text{Sappiamo: } span\{\beta\}=\mathbb{R}^n \\
    \Rightarrow dim(Im(F))=dim(Im(F(\vec{v_i})))+dim(Im(F(\vec{u_i}))) \\
    \to dim(Im(F))=dim(Im(F(\vec{u_i}))) \quad (=n-k)\\  
    \to Im(F)=n-k \\
    \text{Quindi: } \\
    dim(\beta)=k-(n-k) \\ 
    \to dim(\beta)=k-n+k \\ 
    \to dim(\beta)=n \\
    \to \top \\ 
\end{align*}

\section{Sistemi lineari}
un sistema lineare e'
\begin{align*}
    [A](\vec{x})=(\vec{b})
\end{align*}
\subsection{Teorema di struttura di sistemi lineari}: le soluzioni di un sistema lineare sono definite da
\begin{align*}
    \text{sol.}=\{\vec{x_i}+\vec{z} \mid \vec{z}\in Ker(A)\}
\end{align*}
con $\vec{x_i}$ uguale ad una delle soluzioni di $\vec{x}$.
\subsection{Teorema di Rouche'-Capelli}
\begin{itemize}
    \item $rk(A)=rk(A\mid \vec{b})=n \to \exists! $ sol. \\
    Dimostrazione:
    \begin{align*}
        \text{Definiamo: } F : \mathbb{R}^n \to \mathbb{R}^n \\
        A\vec{x}=\vec{b}\\
        \to F(\vec{x})=\vec{b} \\
        \to F^{-1}(\vec{b})=\vec{x} \\
        \text{dato che e' invertibile vuol dire che e' anche biettiva, quindi: } \exists! \, \vec{x}
    \end{align*}
    \item $rk(A)=rk(A\mid \vec{b}) < n \to \infty $ sol. dipendenti da $n-rk(A)$ parametri (o incognite). \\
    Dimostrazione: 
    \begin{align*}
            \text{Definiamo: } F : \mathbb{R}^n \to \mathbb{R}^m \\
            F(\vec{x})=\vec{b}\\
            \text{Dato che $\vec{b}\neq\{\vec{0_v}\}$ allora F non sara' iniettiva} \\
            \exists \, \vec{x}=\vec{b}
    \end{align*}
\end{itemize}

\section{Determinante e inversa di Matrice}
Il determinante e' una funzione unica di ogni matrice quadrata. \\
\vspace{1em} \\
\textbf{Teorema di Binet}: $det(AB)=det(A)det(B)$
\vspace{1em}
\subsection{Teorema del criterio di invertibilità}
\[
\det(A) \neq 0 \iff \exists A^{-1} \text{ tale che } A A^{-1} = I
\]

\textbf{Dimostrazione:}
\begin{itemize}
    \item[$\Rightarrow$] Se $\det(A) \neq 0$, allora $A$ è invertibile. Infatti, l'algoritmo di Gauss-Jordan può trasformare $A$ nella matrice identità, quindi esiste $A^{-1}$.
    \item[$\Leftarrow$] Se esiste $A^{-1}$, allora esiste una matrice tale che $A A^{-1} = I$. Poiché $\det(I) = 1$, allora:
    \[
    \det(A A^{-1}) = \det(A) \cdot \det(A^{-1}) = 1 \Rightarrow \det(A) \neq 0
    \]
\end{itemize}
\vspace{1em} 
\subsection{Teorema fondamentale dell'algebra lineare}
Ricapitolando per $F : \mathbb{R}^n \to \mathbb{R}^n$:
\begin{itemize}
    \item $det(A)\neq0$
    \item $\exists \, A^{-1}$
    \item $F \text{ e' isomorfismo}$
    \item $F\text{ e' biettiva}$
    \item $Ker=\{\vec{0_v}\}$
    \item $dim(Im)=n$
    \item $A\vec{x}=\vec{b} \to \exists! \, \text{sol.}$
    \item $rk(A)=n$
    \item \text{tutti i vettori riga e colonna sono lin. indip.}
\end{itemize}
Dimostrazione:
\begin{align*}
    \text{Ipotizzo: } det(A)\neq0 \\
    \iff \exists! \, F^{-1}\\
    \iff (\text{essendo invertibile e' anche isomorfismo}): span\{[A]\} \equiv_n \mathbb{R}^n \\
    \iff \text{F e' biettiva}\\
    \iff \text{Essendo suriettiva: } Ker=\{\vec{0_v}\} \\
    \iff \text{Essendo iniettiva: } \exists! \, \vec{x} \to A\vec{x}=\vec{b} \to dim(Im)=\mathbb{R}^n \\
    \iff rr(A)=rc(A)=n \to rk(A)=n \\
    \iff rr(A) \to \text{vettori riga indipendenti} \\
    \iff rc(A) \to \text{vettori colonna indipendenti} \\
\end{align*}
\vspace{15em}\\
\section{Cambio di base}
Sia \( F: V \to W \) un'applicazione lineare, con basi \( \beta = \{b_1, \dots, b_n\} \) di \( V \) e \( \beta' = \{b'_1, \dots, b'_m\} \) di \( W \). Allora:
\[
[F(\vec{v})]_{\beta'} = [A]_{\beta, \beta'} [\vec{v}]_{\beta}
\]
dove:
\[
[A]_{\beta, \beta'} = 
\begin{bmatrix}
[F(b_1)]_{\beta'} & [F(b_2)]_{\beta'} & \cdots & [F(b_n)]_{\beta'}
\end{bmatrix}
\]
(con le colonne date dalle coordinate in \( \beta' \) delle immagini dei vettori della base \( \beta \)).
\vspace{1em}\\
\textbf{Caso particolare: matrice identita'}\\
Dato $F : \mathbb{R}^n \to \mathbb{R}^n$
\begin{align*}
    \beta=\{\vec{v_1},\dots,\vec{v_n}\}, \quad \beta'=\{\vec{w_1},\dots,\vec{w_n}\} \\
[A]_{\beta, \beta'} = 
\begin{bmatrix}
[F(\vec{v_1})]_{\beta'} & [F(\vec{v_2})]_{\beta'} & \cdots & [F(\vec{v_n})]_{\beta'}
\end{bmatrix} \\
\to \begin{bmatrix}
(\vec{w_1})_{\beta'} & (\vec{w_2})_{\beta'} & \cdots & (\vec{w_n})_{\beta'}
\end{bmatrix} \\
\to [I_n]
\end{align*}
Quindi la matrice della applicazione lineare associata ad due basi che generano il dominio ed codominio sara' quella di identita'.

\section{Autovalori e autovettori}
Si chiama autovettore un vettore $\vec{v}$ di autovalore $\lambda_i$ della applicazione linerare F se:
\begin{align*}
    F(\vec{v})=\lambda\vec{v}, \quad \forall \lambda_i \in \mathbb{Z}
\end{align*}
con $\mathbb{Z}$ come campo degli scalari; calcolati con
\begin{align*}
    det(A-\lambda_iI)=0
\end{align*}
\vspace{1em}\\
Un'applicazione lineare \( F: V \to V \) è \textbf{diagonalizzabile} se esiste una base \( \beta \) di \( V \) formata da autovettori di \( F \).

In altre parole, esiste una base \( \beta = \{v_1, \dots, v_n\} \) tale che
\[
F(v_i) = \lambda_i v_i, \quad \lambda_i \in \mathbb{K}, \quad i=1,\dots,n,
\]
dove \( \lambda_i \) sono gli autovalori di \( F \).

Allora la matrice di \( F \) rispetto alla base \( \beta \) è diagonale:
\[
[A]_{\beta} = 
\begin{bmatrix}
\lambda_1 & 0 & \cdots & 0 \\
0 & \lambda_2 & \cdots & 0 \\
\vdots & \vdots & \ddots & \vdots \\
0 & 0 & \cdots & \lambda_n
\end{bmatrix}.
\]

Equivalentemente, una matrice \( A \) è diagonalizzabile se esiste una matrice invertibile \( P \) tale che
\[
P^{-1} A P = D,
\]
dove \( D \) è una matrice diagonale. \\
\vspace{1em} \\
\textbf{Teorema}: L'applicazione lineare \( T: \mathbb{R}^n \to \mathbb{R}^n \) è diagonalizzabile se e solo se esiste una base di \( \mathbb{R}^n \) formata da autovettori di \( T \). \\
\vspace{5px} \\
Dimostrazione: 
\begin{itemize}
    \item[\(\Rightarrow\)] Se \(T\) è diagonalizzabile \(\to \exists \beta=\{\vec{a}_1,\dots,\vec{a}_n\}\), e quindi questa base genera \(\mathbb{R}^n\), quindi \(\beta \subseteq \mathbb{R}^n\)
    \item[\(\Leftarrow\)] Se \(\exists \beta=\{\vec{a}_1,\dots,\vec{a}_n\}\) per definizione sappiamo che allora \(T\) è diagonalizzabile.
\end{itemize}
\vspace{1em}
Si definisce \textbf{polinomio caratteristico} di $F$
\begin{align*}
    p_{F}(\lambda)=det(A-\lambda I) \\
    \text{ovvero} \\
    p_{F}(\lambda_i)=0 \\ 
\end{align*}
\vspace{5px} \\
Due matrici $A$ e $B$ si dicono \textbf{simili} se 
\begin{align*}
    B=P^{-1}AP
\end{align*}
Notiamo come anche la matrice diagonalizzabile $D$ sia simile ad A. \\
\vspace{3px}\\
\textbf{Teorema del polinomio caratteristico comune fra matrici simili}: se $A$ e $B$ sono matrici simili allora avranno come polinio caratteristico entrambi lo stesso. \\
\vspace{3px} \\
Dimostrazione: \\
\begin{align*}
    \text{Dato: } B=P^{-1}AP \\
    \Rightarrow det(B)=det(P^{-1}AP) \\
    \Rightarrow det(B)=det(P^{-1})det(A)det(P) \\
    \to det(B)=det(A)det(P^{-1}P) \\
    \to det(B)=det(A) \\
    \to det(B -\lambda I)=det(A -\lambda I) \\
\end{align*}
\subsection{Autospazi}
Si dice autospazio $V$
\begin{align*}
    V=span\{\vec{v} \mid F(\vec{v})=\lambda \vec{v}\} \\
    = Ker(A-\lambda I)
\end{align*}
\textbf{Teorema:} Gli autovettori di autovalori diversi tra loro sono linearmente indipendenti.
\begin{align*}
    \lambda_1\neq\dots\neq\lambda_n\\
    \to \text{Gli autovettori $\vec{a_1},\dots,\vec{a_n}$ sono linearmente indipendenti}
\end{align*}
\vspace{1em} \\
Indichiamo con \textbf{molteplicita' algebrica} di $\lambda_i$ il grado del polinio $p_{F}(\lambda_i)$. Questa  se coincede con la \textbf{molteplicita' geometrica}, ovvero la dimensione di $Ker(A-\lambda_i I)$, comporta che la matrice $A$ e' diagonalizzabile.

\section{Prodotto scalare}
Due vettori sono ortogonali se $\langle \vec{u_1},\vec{u_2}\rangle=0$. \\
\textbf{Base ortogonale}: e' una base di vettori ortogonali
\begin{align*}
    \beta_{W^{\bot}}=\{\vec{u_1},\vec{u_2}\}
\end{align*}
\textbf{Base ortonormale}: e' una base composta da vettori normalizzati ortogonali
\begin{align*}
    \mid\mid \vec{e_i} \mid\mid =1 \\
    \beta_{ortonormale}=\{\vec{e_1},\vec{e_2}\} \\
\end{align*}
\subsection{Processo di Gram-Schmidt}
Serve per calcolare una base ortonormale (o ortogonale) di un sottospazio vettoriale.
\begin{align*}
    W \subseteq \mathbb{R}^n \\
    \beta_W=\{\vec{v_1},\vec{v_2},\dots,\vec{v_n}\} \\
    \beta_{W^{\bot}}=\beta_{ortogonale}=\{\vec{u_1},\vec{u_2},\dots,\vec{u_n}\} \\
    \beta_{ortonormale}=\{\vec{e_1},\vec{e_2},\dots,\vec{e_n}\} \\
    \vec{u_1}=\vec{v_1}, \quad \vec{e_1}=\frac{\vec{u_1}}{|| \vec{u_1} ||} \\
    \vec{u_i}=\vec{v_i}-\sum_{j=1}^{i-1} \text{proj}_{\vec{u_j}}(v_i), \quad \vec{e_i}=\frac{\vec{u_i}}{||\vec{u_i}||}
\end{align*}
\end{document} 